\SixthArticle{18}
We now introduce the following notation:
\[
[38]\qquad \Pi(k,z)=\frac{1\mathbin{.}2\mathbin{.}3\mathinner{.\,.\,.\,.\,.}k}{(z+1)(z+2)(z+3)\mathinner{.\,.\,.\,.}(z+k)}k^z
\]
where $k$ is naturally understood to denote a positive integer; with this restriction, $\Pi(k,z)$ represents a completely determined function of the two quantities $k,z$. It is then easy to see that the theorem proposed at the end of the preceding article can be written as follows:
\[
[39]\qquad F(\alpha,\beta,\gamma,1)=\frac{\Pi(k,\gamma-1)\mathbin{.}\Pi(k,\gamma-\alpha-\beta-1)}{\Pi(k,\gamma-\alpha-1)\mathbin{.}\Pi(k,\gamma-\beta-1)}\mathbin{.}F(\alpha,\beta,\gamma+k,1)
\]
\SixthArticle{19}
It will be worth examining the nature of the function $\Pi(k,z)$ more closely. Whenever $z$ is a negative integer, the function evidently takes an infinitely large value as soon as a sufficiently large value is assigned to $k$. For nonnegative integer values of $z$, however, we have
\[
\begin{aligned}
\Pi(k,0)&=1\\[5pt]
\Pi(k,1)&=\frac{1}{1+\frac1k}\\[5pt]
\Pi(k,2)&=\frac{1\mathbin{.}2}{(1+\frac1k)(1+\frac2k)}\\[5pt]
\Pi(k,3)&=\frac{1\mathbin{.}2\mathbin{.}3}{(1+\frac1k)(1+\frac2k)(1+\frac3k)}
\end{aligned}
\]
etc., or in general
\[
[40]\qquad \Pi(k,z)=\frac{1\mathbin{.}2\mathbin{.}3\mathinner{.\,.\,.\,.\,.\,.}z}{(1+\frac1k)(1+\frac2k)(1+\frac3k)\mathinner{.\,.\,.\,.}(1+\frac zk)}
\]
\SourceLeaf{154}
In general, however, for \textit{any} value of $z$ we have
\[
[41]\qquad \Pi(k,z+1)=\Pi(k,z)\mathbin{.}\frac{1+z}{1+\frac{1+z}{k}}
\]
\[
[42]\qquad \Pi(k+1,z)=\Pi(k,z)\mathbin{.}\left\{\frac{(1+\frac1k)^{z+1}}{1+\frac{1+z}{k}}\right\}
\]
and hence, since $\Pi(1,z)=\frac1{z+1}$,
\[
[43]\qquad \Pi(k,z)=\frac1{z+1}\mathbin{.}\frac{2^{z+1}}{1^z\mathbin{.}(2+z)}\mathbin{.}\frac{3^{z+1}}{2^z(3+z)}\mathbin{.}\frac{4^{z+1}}{3^z(4+z)}\mathinner{.\,.\,.\,.}\frac{k^{z+1}}{(k-1)^z(k+z)}
\]
\SixthArticle{20}
The \textit{limit} towards which, for a given value of $z$, the function $\Pi(k,z)$ steadily converges as $k$ grows to infinity deserves particular attention. First let $h$ be a finite value of $k$ greater than $z$. It is then clear that, if $k$ is supposed to pass from $h$ to $h+1$, the logarithm of $\Pi(k,z)$ receives an increment expressed by the following convergent series:
\[
\frac{z(1+z)}{2(h+1)^2}+\frac{z(1-zz)}{3(h+1)^3}+\frac{z(1+z^3)}{4(h+1)^4}+\frac{z(1-z^4)}{5(h+1)^5}+\quad\text{etc.}
\]
Thus, if $k$ passes from the value $h$ to $h+n$, the logarithm of $\Pi(k,z)$ will receive the increment
\[
\begin{aligned}
&\tfrac12 z(1+z)\left(\frac1{(h+1)^2}+\frac1{(h+2)^2}+\frac1{(h+3)^2}+\text{etc.}+\frac1{(h+n)^2}\right)\\[5pt]
&+\tfrac13 z(1-zz)\left(\frac1{(h+1)^3}+\frac1{(h+2)^3}+\frac1{(h+3)^3}+\text{etc.}+\frac1{(h+n)^3}\right)\\[5pt]
&+\tfrac14 z(1+z^3)\left(\frac1{(h+1)^4}+\frac1{(h+2)^4}+\frac1{(h+3)^4}+\text{etc.}+\frac1{(h+n)^4}\right)\\[5pt]
&+\text{etc.}
\end{aligned}
\]
which can readily be shown to remain finite even if $n$ grows to infinity. Thus, unless an infinite factor was already present in $\Pi(h,z)$, that is, unless $z$ is a negative integer, the limit of $\Pi(k,z)$ for $k=\infty$ will certainly be a finite quantity. Hence $\Pi(\infty,z)$ evidently depends only on $z$, or represents a fully determined function of $z$, which from now on will be denoted simply by $\Pi z$. We therefore define the function $\Pi z$ by the value of the product
\[
\frac{1\mathbin{.}2\mathbin{.}3\mathinner{.\,.\,.\,.\,.}k\mathbin{.}k^z}{(z+1)(z+2)(z+3)\mathinner{.\,.\,.\,.}(z+k)}
\]
\SourceLeaf{155}
for $k=\infty$, or, if you prefer, by the limit of the infinite product
\[
\frac{1}{z+1}.\frac{2^{z+1}}{1^z(2+z)}.
\frac{3^{z+1}}{2^z(3+z)}.\frac{4^{z+1}}{3^z(4+z)}\quad\mathrm{etc.}
\]

\Needspace{24\baselineskip}
\SixthArticle{21}
Equation 41 immediately gives the fundamental equation
\[
\begin{array}{@{}ll@{}}
[44]&\displaystyle\Pi(z+1)=(z+1)\Pi z
\end{array}
\]
and hence in general, with $n$ denoting any positive integer,
\[
\begin{array}{@{}ll@{}}
[45]&\displaystyle\Pi(z+n)=(z+1)(z+2)(z+3)\,.\,.\,.\,.(z+n)\Pi z
\end{array}
\]
For a negative integer value of $z$ the value of the function $\Pi z$ will be infinitely large; for nonnegative integer values we have
\[
\begin{array}{rcl}
\Pi0&=&1\\
\Pi1&=&1\\
\Pi2&=&2\\
\Pi3&=&6\\
\Pi4&=&24\quad\mathrm{etc.}
\end{array}
\]
and in general
\[
\begin{array}{@{}ll@{}}
[46]&\displaystyle\Pi z=1.2.3\,.\,.\,.\,.z
\end{array}
\]
But this property of our function is wrongly passed off as its definition: it is by its nature restricted to integer values, and is shared by infinitely many other functions besides ours (for example, $\cos2\pi z.\Pi z$, $\cos\pi z^{2n}\Pi z$ etc., where $\pi$ denotes half the circumference of a circle whose radius $=1$).

